% Originally: content/week-07.tex, \section{Solutions}

\section{Solutions}
\label{sec:submanifolds_solutions}

% Transition: Claude Opus 5 (High)
Not every solution in this chapter is collected here. \Cref{ex:dimension_submanifold_system}
(\cpageref{ex:dimension_submanifold_system}) and \cref{ex:submanifold_matrices_rank1}
(\cpageref{ex:submanifold_matrices_rank1}) are solved where they appear.

\begin{ainote}
Corsin's notes for this material cover only the theory (submanifolds, the regular value and local
parametrization theorems, tangent and normal spaces), with no page working out any of the
official-sheet exercises, so all solutions below are original. Those answering a problem from
Sheet 7 have been checked against \texttt{exercises/Sol7\_Analysis2\_eng.pdf}, and no divergence
from the official solutions was found. The exam problems mined into this chapter come from papers
that ship no key at all, so those solutions rest on independent derivation alone.
\end{ainote}

\begin{exercisesolution}[Solution to \cref{ex:7.6}]
% Generator: Claude Sonnet 5 (Medium)
\begin{enumerate}[label=\textbf{(\alph*)}]
  \item Fixing $\theta_0,\varphi_0$, the image of $r \mapsto \Phi(r,\theta_0,\varphi_0)$ is the
    radial half-line through the point $p_0 := (\sin\theta_0\cos\varphi_0,
    \sin\theta_0\sin\varphi_0, \cos\theta_0)$ on the unit sphere, i.e., $\{rp_0 : r \in
    (0,\infty)\}$. Fixing $r_0,\varphi_0$, the image of $\theta \mapsto \Phi(r_0,\theta,\varphi_0)$
    is a semicircle of radius $r_0$ from the north pole to the south pole (excluding the poles),
    at longitude $\varphi_0$. Fixing $r_0,\theta_0$, the image of $\varphi \mapsto
    \Phi(r_0,\theta_0,\varphi)$ is a circle of latitude of radius $r_0\sin\theta_0$ at height
    $r_0\cos\theta_0$ (excluding the point at $\varphi = \pm\pi$).
  \item For each fixed $r$, the image sweeps out the sphere of radius $r$ minus the half-plane
    at $\varphi = \pm\pi$ (the \emph{anti}-meridian, since the prime meridian is $\varphi = 0$).
    Overall,
    \[\operatorname{im}(\Phi) = \mathbb{R}^3\setminus\big((-\infty,0]\times\{0\}\times\mathbb{R}\big).\]
  \item Computing,
    \[D_{(r,\theta,\varphi)}\Phi = \begin{pmatrix} \sin\theta\cos\varphi & r\cos\theta\cos\varphi & -r\sin\theta\sin\varphi \\ \sin\theta\sin\varphi & r\cos\theta\sin\varphi & r\sin\theta\cos\varphi \\ \cos\theta & -r\sin\theta & 0 \end{pmatrix},\]
    hence, expanding along the last row,
    \[\begin{aligned}
    \det(D_{(r,\theta,\varphi)}\Phi) &= r^2\sin\theta\cos^2\theta\cos^2\varphi+r^2\sin^3\theta\sin^2\varphi+r^2\sin\theta\cos^2\theta\sin^2\varphi+r^2\sin^3\theta\cos^2\varphi \\
    &= r^2\big(\sin\theta\cos^2\theta+\sin^3\theta\big) = r^2\sin\theta.
    \end{aligned}\]
  \item Since $r \in (0,\infty)$ and $\theta \in (0,\pi)$, $r^2\sin\theta$ is never $0$, so
    $\Jac\Phi(r,\theta,\varphi)$ is everywhere invertible and $\Phi$ is a local diffeomorphism by
    \Cref{thm:inverse_function_theorem}. To see it is a \textbf{global} diffeomorphism onto its
    image, it remains to check $\Phi$ is injective: if $\Phi(r_1,\theta_1,\varphi_1) =
    \Phi(r_2,\theta_2,\varphi_2)$, then $r_1=r_2$ follows by taking the norm of both sides; since
    $\cos : (0,\pi) \to (-1,1)$ is injective, the third coordinate gives $\theta_1=\theta_2$;
    finally $(\cos\varphi,\sin\varphi)$ with $\varphi \in (-\pi,\pi)$ determines a unique point on
    the unit circle, so $\varphi_1=\varphi_2$. Hence $\Phi$ is injective, and (being a local
    diffeomorphism everywhere on its domain) a diffeomorphism onto its image.
\end{enumerate}
\end{exercisesolution}

\begin{exercisesolution}[Proof of \cref{ex:submanifold_quadric_tangent_space}]
% Generator: Gemini 3.6 Flash (High)
% Correction: Claude Opus 5 (High) --- ^T -> \transp throughout, \implies replaced by prose, and the
% check that p lies on M added: the tangent space computation presupposes it.
\begin{enumerate}[label=\textbf{(\alph*)}]
  \item Define the smooth function $F : \mathbb{R}^3 \to \mathbb{R}$,
    $F(x_1,x_2,x_3) := x_1 x_2 + x_2 x_3 + x_1 x_3 - 1$, so that $M = F^{-1}(0)$. Its gradient is
\[\nabla F(x) = (x_2+x_3,\; x_1+x_3,\; x_1+x_2)\transp .\]
    Suppose $\nabla F(x) = 0$. Adding the three equations $x_2+x_3 = x_1+x_3 = x_1+x_2 = 0$ gives
    $2(x_1+x_2+x_3) = 0$, and subtracting each equation in turn then forces $x_1 = x_2 = x_3 = 0$.
    But $F(0,0,0) = -1 \neq 0$, so the origin is not in $M$. Hence $\nabla F$ is non-zero at every
    point of $M$, which is to say $0$ is a regular value of $F$. By the regular value theorem
    (\cref{thm:regular_value_theorem}), $M$ is a smooth submanifold of $\mathbb{R}^3$ of dimension
    $3 - 1 = 2$.

  \item First, $p = (-1,2,3)$ does lie on $M$: indeed $(-1)(2) + (2)(3) + (-1)(3) = -2 + 6 - 3 = 1$.
    At this point,
\[\nabla F(-1,2,3) = (2+3,\; -1+3,\; -1+2)\transp = (5, 2, 1)\transp .\]
    The tangent space is the kernel of $DF(p)$, equivalently the orthogonal complement of
    $\nabla F(p)$:
\[T_p M = \left\{(v_1, v_2, v_3)\transp \in \mathbb{R}^3 \;\middle|\; 5v_1 + 2v_2 + v_3 = 0\right\}.\]
    Reading $v_3$ off from $v_1$ and $v_2$ gives two independent solutions: taking $(v_1,v_2) =
    (1,0)$ yields $v_3 = -5$, and taking $(v_1,v_2) = (0,1)$ yields $v_3 = -2$. Hence
\[\left\{\begin{pmatrix} 1 \\ 0 \\ -5 \end{pmatrix}, \begin{pmatrix} 0 \\ 1 \\ -2 \end{pmatrix}\right\}\]
    is a basis of $T_p M$; the two vectors are independent because their first two components
    already are, and there are two of them, matching $\dim T_p M = 2$.
\end{enumerate}
\begin{ainote}
The official solution gives the basis $\left\{(2,-5,0)\transp,\ (1,0,-5)\transp\right\}$, which
looks different from the one above but is not a disagreement: both are bases of the same plane
$\ker(5,2,1)$, and each vector of either list satisfies $5v_1+2v_2+v_3 = 0$. \qt{Give a basis} has
infinitely many correct answers, and no canonical one --- so when checking work against a master
solution here, verify that each vector is annihilated by $\nabla F(p)$ and that the two are
independent, rather than comparing entries. The two lists even share the vector $(1,0,-5)\transp$.
The official solution also invokes the constant rank theorem via surjectivity of $D_xF$ where we
cite the regular value theorem (\cref{thm:regular_value_theorem}); for a real-valued $F$ these say
the same thing, since a linear map $\mathbb{R}^3 \to \mathbb{R}$ is surjective precisely when its
gradient is non-zero.
\end{ainote}
\end{exercisesolution}

\begin{exercisesolution}[Solution of \cref{ex:image_of_segment_submanifold}]
% Generator: Claude Opus 5 (High)
For $\alpha = 0$ the determinant computed in \cref{ex:local_diffeo_parameter_alpha} is
$-16 \neq 0$, so fix $r > 0$ small enough that $\Phi$ restricted to $B_r(0,0)$ is a
$C^1$-diffeomorphism onto the open set $\Phi(B_r(0,0))$. Write
\[S := \{(x,y) \in B_r(0,0) \mid y = 0\} = (-r,r) \times \{0\}\]
for the segment in question. It is itself a $1$-dimensional $C^1$ submanifold of $\mathbb{R}^2$:
given $p = (a,0) \in S$, the translation $\Psi(x,y) := (x-a,\, y)$ is a diffeomorphism from
$U_p := B_r(0,0)$ onto the open set $V := \Psi(U_p)$, it sends $p$ to $0$, and
$\Psi(U_p \cap S) = (\mathbb{R}\times\{0\}) \cap V$.
\begin{enumerate}[label=\textbf{(\alph*)}]
  \setcounter{enumi}{3}
  \item \textbf{True.} Apply \cref{prop:diffeomorphism_preserves_submanifold} to the
    $C^1$-diffeomorphism $\Phi|_{B_r(0,0)} : B_r(0,0) \to \Phi(B_r(0,0))$ and to $S$. The image
    $\Phi(S)$ is a $C^1$ submanifold of $\mathbb{R}^2$ of dimension $1$.
  \item \textbf{False.} By \textbf{(d)} the set is $1$-dimensional, and no non-empty set is a
    submanifold of two different dimensions at once. Independently of that: a $2$-dimensional
    submanifold of $\mathbb{R}^2$ is an \emph{open} subset of $\mathbb{R}^2$, because a chart with
    $m = n = 2$ reads $\Psi(U \cap M) = V = \Psi(U)$, so $U \cap M = U$ and $M$ contains the open
    neighbourhood $U$ of each of its points. But $\Phi(S)$ is not open. Its complement in
    $\Phi(B_r(0,0))$ is the image of $B_r(0,0) \setminus S$, which is dense in $B_r(0,0)$, and a
    homeomorphism carries dense sets to dense sets, so every point of $\Phi(S)$ is a limit of
    points outside it.
\end{enumerate}
\begin{remark}[The pair is a dimension check, not a submanifold check]
The two items ask about the same set, so at most one can be true, and the real question is which
number a diffeomorphism is allowed to change. It changes none:
\cref{rem:dimension_is_a_diffeomorphism_invariant} is the whole content of the block. The picture
that makes \textbf{(e)} tempting is that $\Phi$ maps a two-dimensional ball onto a two-dimensional
region, so its image \qt{fills out} the plane. It does, but the segment does not, and it is the
segment being transported.
\end{remark}
\end{exercisesolution}

% Generator: Claude Opus 5 (High)
% No solution key ships with the Jossen papers, so this is an independent derivation.
\begin{exercisesolution}[\cref{ex:orthogonal_group_submanifold}]
\begin{enumerate}[label=\textbf{(\alph*)}]
  \item With $A = \left(\begin{smallmatrix} a & b \\ c & d\end{smallmatrix}\right)$,
\[AA\transp = \begin{pmatrix} a & b \\ c & d \end{pmatrix}
              \begin{pmatrix} a & c \\ b & d \end{pmatrix}
            = \begin{pmatrix} a^2+b^2 & ac+bd \\ ac+bd & c^2+d^2 \end{pmatrix},\]
    so in the coordinates of the statement
\[F(a,b,c,d) = \big(a^2+b^2,\ ac+bd,\ ac+bd,\ c^2+d^2\big).\]
    Every component is a polynomial, hence $F$ is smooth.

  \item $(AA\transp)\transp = (A\transp)\transp A\transp = AA\transp$, so $F(A)$ is symmetric for
    every $A$. This is already visible in the coordinate formula: the second and third components
    of $F$ are the same function.

  \item Differentiating the four components with respect to $a,b,c,d$ in turn,
\[\Jac F(a,b,c,d) = \begin{pmatrix}
  2a & 2b & 0 & 0 \\
  c & d & a & b \\
  c & d & a & b \\
  0 & 0 & 2c & 2d
\end{pmatrix}.\]

  \item The second and third rows coincide for every $A$, so $\rank \Jac F(A) \le 3$ always. To see
    that equality holds on $\Orth_2(\mathbb{R})$, note that $AA\transp = \id$ says exactly that the
    two rows $u := (a,b)$ and $v := (c,d)$ of $A$ are orthonormal. Suppose
\[\alpha(2a,2b,0,0) + \beta(c,d,a,b) + \gamma(0,0,2c,2d) = 0.\]
    Its first two coordinates read $2\alpha u + \beta v = 0$, and $u,v$ are linearly independent,
    so $\alpha = \beta = 0$. The last two then read $2\gamma v = 0$, so $\gamma = 0$ as well. The
    three distinct rows are therefore independent and the rank is exactly $3$.

    This is why \cref{thm:regular_value_theorem} cannot be applied to $F$ as a map
    $\mathbb{R}^4 \to \mathbb{R}^4$: the theorem needs $D_AF$ to be \emph{surjective}, that is of
    rank $4$, and the rank is $3$ at every point of $\Orth_2(\mathbb{R})$. Nor is this a defect of
    the particular level: by \textbf{(b)} the image of $F$ lies in a three-dimensional subspace, so
    $D_AF$ is nowhere surjective, for any $A$ whatsoever.

  \item Identify $\operatorname{Sym}_2(\mathbb{R}) \cong \mathbb{R}^3$ by
    $\left(\begin{smallmatrix} p & q \\ q & s\end{smallmatrix}\right) \mapsto (p,q,s)$ and let
\[\widetilde F : \mathbb{R}^4 \to \mathbb{R}^3, \qquad
  \widetilde F(a,b,c,d) := \big(a^2+b^2,\ ac+bd,\ c^2+d^2\big),\]
    which is $F$ with the duplicated component dropped. Then
    $\Orth_2(\mathbb{R}) = \widetilde F^{-1}\{(1,0,1)\}$, and
\[\Jac \widetilde F(a,b,c,d) = \begin{pmatrix}
  2a & 2b & 0 & 0 \\
  c & d & a & b \\
  0 & 0 & 2c & 2d
\end{pmatrix}\]
    has rank $3$ at every point of $\Orth_2(\mathbb{R})$, by the computation in \textbf{(d)}. So
    $(1,0,1)$ is a regular value of $\widetilde F$, and the regular value theorem makes
    $\Orth_2(\mathbb{R})$ a smooth submanifold of $\mathbb{R}^4 = \operatorname{Mat}_2(\mathbb{R})$
    of dimension $4 - 3 = 1$.

  \item By \cref{prop:gradient_orthogonality_level_sets} applied to each of the three defining
    equations, or directly from the description of the tangent space to a regular level set,
    $T_{\id}\Orth_2(\mathbb{R}) = \ker \Jac \widetilde F(\id)$. At
    $(a,b,c,d) = (1,0,0,1)$,
\[\Jac \widetilde F(1,0,0,1) = \begin{pmatrix}
  2 & 0 & 0 & 0 \\
  0 & 1 & 1 & 0 \\
  0 & 0 & 0 & 2
\end{pmatrix},\]
    so a vector $(v_1,v_2,v_3,v_4)$ lies in the kernel precisely when $v_1 = 0$, $v_4 = 0$ and
    $v_2 + v_3 = 0$. Hence
\[T_{\id}\Orth_2(\mathbb{R}) = \operatorname{span}\left\{
  \begin{pmatrix} 0 \\ 1 \\ -1 \\ 0\end{pmatrix}\right\}
  = \operatorname{span}\left\{\begin{pmatrix} 0 & 1 \\ -1 & 0\end{pmatrix}\right\},\]
    a line, in agreement with $\dim \Orth_2(\mathbb{R}) = 1$ from \textbf{(e)}. Using the full
    $4\times4$ Jacobian of \textbf{(c)} gives the same kernel, since the repeated row imposes no
    new condition.
\end{enumerate}
\begin{remark}[The tangent space names the answer]
The line found in \textbf{(f)} consists of the \emph{antisymmetric} $2\times2$ matrices, and that
is not a coincidence: differentiating $A(t)A(t)\transp = \id$ along any curve in
$\Orth_2(\mathbb{R})$ through $\id$ gives $\dot A(0) + \dot A(0)\transp = 0$. The same computation
in $n$ dimensions gives $\dim\Orth_n(\mathbb{R}) = \binom{n}{2}$, the dimension of the
antisymmetric matrices, which for $n=2$ is $1$.

That $\Orth_2(\mathbb{R})$ is one-dimensional is worth pausing on, because it is easy to expect
more. Writing $u = (\cos\theta,\sin\theta)$ for the first row, orthonormality leaves exactly two
choices of second row, so $\Orth_2(\mathbb{R})$ is a disjoint union of two circles: the rotations
$\det A = 1$ and the reflections $\det A = -1$. A submanifold need not be connected, and the
regular value theorem never claims it is.

Set against \cref{ex:submanifold_matrices_rank1}, which lives in the same $\mathbb{R}^4$, the
contrast is instructive. There a single scalar equation cut out a three-dimensional submanifold and
the regular value theorem applied as it stands. Here the naive count would suggest
$4 - 4 = 0$ dimensions, and it is wrong precisely because the four equations $AA\transp = \id$ are
not independent: one of them is a duplicate. Counting equations is only a shortcut for computing a
rank, and when the shortcut and the rank disagree it is the rank that is right.
\end{remark}
\end{exercisesolution}

\begin{exercisesolution}[Solution to \cref{ex:operations_preserving_submanifolds}]
% Generator: Claude Opus 5 (High) --- the four statements are the examiner's; the working, the
% counterexample in (b) and the remark are written for these notes. Probeprfg1 ships no key.
Throughout, \qt{submanifold} is used in the sense of \cref{def:submanifold}: locally
straightenable by a diffeomorphism of the ambient space.
\begin{enumerate}[label=\textbf{(\alph*)}]
  \item \textbf{True}, and there is nothing to prove beyond intersecting with one more open set.
    Being open in $M$ means $N = M \cap W$ for some open $W \subseteq \mathbb{R}^n$. Fix $p \in
    N$ and take a submanifold chart $\Psi : U_p \to V$ for $M$ at $p$. Replacing $U_p$ by
    $U_p \cap W$, which is still an open neighbourhood of $p$, turns it into a submanifold chart
    for $N$ at $p$ of the same dimension. The definition is local, so shrinking the neighbourhood
    costs nothing.
  \item \textbf{False}, and this is the one that has to be met rather than reasoned about.
    In $\mathbb{R}^2$ take
    \[M := \{(x,0) \mid x > 0\}, \qquad N := \{(0,y) \mid y \in \mathbb{R}\}.\]
    They are disjoint, since every point of $M$ has $x > 0$ and every point of $N$ has $x = 0$,
    and each is a $1$-dimensional submanifold: $M$ is an open subset of the $x$-axis, so
    part \textbf{(a)} applies, and $N$ is a line.

    But $M \cup N$ is not a submanifold at the origin. Every ball $B_\varepsilon(0,0)$ meets it
    in a segment of the $y$-axis together with a half-open segment running right, and deleting
    the origin from that set leaves \emph{three} connected pieces, whereas deleting a point from
    a straightened $(\mathbb{R}\times\{0\})\cap V$ leaves two. No diffeomorphism can change the
    number of pieces, so no submanifold chart exists at $(0,0)$.

    What goes wrong is that disjointness is not enough: the two pieces are disjoint but one of
    them accumulates on the other. Requiring instead that $M$ and $N$ have disjoint
    \emph{closures} makes the statement true, and then each point of the union has a
    neighbourhood meeting only one of the two.
  \item \textbf{True}, and it is the chart definition again. Fix $(p,q) \in M \times N$ and take
    submanifold charts $\Psi : U_p \to V$ for $M$ and $\Xi : U_q \to W$ for $N$, with
    $\Psi(U_p \cap M) = (\mathbb{R}^k\times\{0\})\cap V$ and
    $\Xi(U_q \cap N) = (\mathbb{R}^l\times\{0\})\cap W$. Then $\Psi\times\Xi$ is a
    diffeomorphism of the open set $U_p\times U_q \subseteq \mathbb{R}^{m+n}$ onto $V\times W$,
    and it carries $(U_p\times U_q) \cap (M\times N)$ onto
    \[\big((\mathbb{R}^k\times\{0\})\cap V\big) \times \big((\mathbb{R}^l\times\{0\})\cap W\big).\]
    Composing with the permutation of coordinates that gathers the $k$ and the $l$ live
    directions together, which is linear and invertible and hence a diffeomorphism, turns this
    into $(\mathbb{R}^{k+l}\times\{0\}) \cap (V\times W)$. That is a submanifold chart of
    dimension $k+l$.
  \item \textbf{True}, and the hypothesis on the tangent spaces is exactly what is needed. Fix
    $p \in M \cap N$. Each of $M$ and $N$ is locally a level set, and \cref{def:submanifold}
    supplies this directly: a submanifold chart $\Psi : U \to V$ for $M$ at $p$ satisfies
    $\Psi(U \cap M) = (\mathbb{R}^{n-1}\times\{0\}) \cap V$, so its last component
    $F := \Psi_n$ is $C^1$ with $M \cap U = F^{-1}\{0\}$, and $\nabla F(p) \neq 0$ because it is
    a row of the invertible matrix $\Jac\Psi(p)$. Do the same for $N$, shrink to a common
    neighbourhood, and write $G$ for its function. So
    \[M \cap U = F^{-1}\{0\}, \quad \nabla F(p) \neq 0, \qquad
      N \cap U = G^{-1}\{0\}, \quad \nabla G(p) \neq 0 .\]
    By \cref{prop:gradient_orthogonality_level_sets}, $T_pM$ has $\nabla F(p)$ as its
    normal line, and likewise for $N$. Two hyperplanes in $\mathbb{R}^n$ coincide precisely when
    their normal lines do, so $T_pM \neq T_pN$ says exactly that $\nabla F(p)$ and $\nabla G(p)$
    are linearly independent.

    That makes $H := (F,G) : U \to \mathbb{R}^2$ a map whose Jacobian at $p$ has the two
    independent rows $\nabla F(p)\transp$ and $\nabla G(p)\transp$, hence rank $2$, so $DH_p$ is
    surjective. Full rank is an open condition, so after shrinking $U$ it holds throughout, and
    \cref{thm:regular_value_theorem} makes
    \[H^{-1}\{0\} = (M\cap U) \cap (N \cap U) = (M\cap N)\cap U\]
    a submanifold of dimension $n-2$. As $p$ was arbitrary, $M \cap N$ is one.
\end{enumerate}

\begin{remark}[Transversality is the name of the hypothesis in \textbf{(d)}]
Two submanifolds are said to meet \newterm{transversally} at $p$ if their tangent spaces together
span the ambient space, $T_pM + T_pN = \mathbb{R}^n$. For two hyperplanes that is the same as
being distinct, which is why part \textbf{(d)} can state the condition the way it does, and the
general theorem says that a transversal intersection is a submanifold of dimension
$\dim M + \dim N - n$. Here $(n-1)+(n-1)-n = n-2$, as found.

The hypothesis is doing real work, and dropping it breaks the conclusion immediately: two spheres
in $\mathbb{R}^3$ meeting at a single point of tangency have the same tangent plane there, and
their intersection is a point, of dimension $0$ rather than $1$. Two coincident planes are worse
still, intersecting in a plane. Transversality is what rules both out, and the reason it is the
right condition rather than a convenient one is visible in the proof: it is precisely the
statement that stacking the two constraints gives a map of full rank.
\end{remark}
\end{exercisesolution}

\begin{exercisesolution}[Solution to \cref{ex:level_set_of_a_wide_jacobian}]
% Generator: Claude Opus 5 (High) --- the six statements are the examiner's; the working is
% written for these notes. Verdicts checked afterwards against exercises_2024/mocksol.pdf and
% agree on all six.
Everything rests on two computations. First, $\gamma(0) = (0,0,1)$, and differentiating
componentwise,
\[\gamma'(t) = \left(2\cos(2t),\ 1+t,\
  \frac{-\sin t\,(1-t) + \cos t}{(1-t)^2}\right), \qquad\text{so}\qquad
  \gamma'(0) = \begin{pmatrix} 2 \\ 1 \\ 1 \end{pmatrix}.\]
Second, $\Jac F(0,0,1)$ has rank $2$: its leftmost $2\times2$ minor is
$\left(\begin{smallmatrix} 1 & 1 \\ 0 & 1\end{smallmatrix}\right)$, of determinant $1$. Full rank
is an open condition, since that minor is a continuous function of the point, so $\Jac F$ has
rank $2$ on a whole neighbourhood $U$ of $(0,0,1)$, and \cref{thm:regular_value_theorem} applied
to $F|_U$ makes $M \cap U$ a submanifold of dimension $3-2 = 1$. Two independent equations cut
two dimensions out of three and leave a curve; parts \textbf{(b)} and \textbf{(c)} follow at
once.
\begin{enumerate}[label=\textbf{(\alph*)}]
  \item \textbf{False.} By the chain rule,
    \[(F\circ\gamma)'(0) = \Jac F(\gamma(0))\,\gamma'(0)
      = \begin{pmatrix} 1 & 1 & -2 \\ 0 & 1 & -1\end{pmatrix}
        \begin{pmatrix} 2 \\ 1 \\ 1 \end{pmatrix}
      = \begin{pmatrix} 2+1-2 \\ 1-1 \end{pmatrix}
      = \begin{pmatrix} 1 \\ 0 \end{pmatrix} \neq \begin{pmatrix} 0 \\ 0\end{pmatrix}.\]
    This is worth reading as information rather than as a failed computation: had $\gamma$ run
    inside $M$, then $F\circ\gamma$ would be constant and the derivative would vanish. So
    $\gamma$ meets $M$ at $t=0$ and immediately leaves it.
  \item \textbf{False:} the dimension is $1$, not $2$.
  \item \textbf{True.}
  \item \textbf{False.} By \cref{def:normal_space} the normal space is
    $N_{(0,0,1)}M = \operatorname{span}\{(1,1,-2)\transp,\ (0,1,-1)\transp\}$, the span of the
    rows of $\Jac F(0,0,1)$. If $\gamma'(0) = a(1,1,-2)\transp + b(0,1,-1)\transp$, the first
    coordinate forces $a = 2$ and the second then forces $b = -1$; but the third coordinate of
    that combination is $-2\cdot2 - (-1) = -3$, not $1$. Note that $\gamma'(0)$ is not tangent
    either: part \textbf{(a)} says exactly that it is not in $\ker\Jac F(0,0,1) = T_{(0,0,1)}M$.
    It is simply a vector in neither subspace, which is what a curve leaving $M$ produces.
  \item \textbf{True.} To solve for $x$ and $z$ in terms of $y$, the block of $\Jac F(0,0,1)$ in
    the $x$- and $z$-columns must be invertible, and
    \[\det\begin{pmatrix} 1 & -2 \\ 0 & -1\end{pmatrix} = -1 \neq 0 .\]
    \Cref{thm:implicit_function_theorem} therefore supplies $\delta > 0$ and $C^1$ functions
    $\varphi,\psi$ on $(-\delta,\delta)$ with $\varphi(0) = 0$, $\psi(0) = 1$ and
    $F(\varphi(y),y,\psi(y)) = (2,-2)\transp$ throughout. This exhibits the curve of part
    \textbf{(c)} concretely, as a graph over the $y$-axis.
  \item \textbf{True.} Complete $F$ to a map of $\mathbb{R}^3$ to itself by appending the
    coordinate that part \textbf{(e)} solved \emph{for}, namely $y$:
    \[\Phi(x,y,z) := \big(F_1(x,y,z),\ F_2(x,y,z),\ y\big), \qquad
      \Jac\Phi(0,0,1) = \begin{pmatrix} 1 & 1 & -2 \\ 0 & 1 & -1 \\ 0 & 1 & 0\end{pmatrix},\]
    whose determinant is $1 \cdot (1\cdot 0 - (-1)\cdot 1) = 1 \neq 0$. By
    \cref{thm:inverse_function_theorem}, $\Phi$ is a $C^1$ diffeomorphism of some
    $B_\varepsilon((0,0,1))$ onto an open subset of $\mathbb{R}^3$. Writing $\pi(u,v,w) := (u,v)$
    for the projection, we have $F = \pi\circ\Phi$, and $\pi$ maps open sets to open sets, since
    $\pi(B_r(p)) = B_r(\pi(p))$. Hence $F(B_\varepsilon) = \pi(\Phi(B_\varepsilon))$ is open.
\end{enumerate}

\begin{ainote}
The examination calls $M$ a \qt{smooth manifold} in parts \textbf{(b)} and \textbf{(c)}, but it
assumes only $F \in C^1$, so what the regular value theorem delivers is a $C^1$ submanifold. The
distinction changes neither verdict, since both parts are about the dimension.
\end{ainote}
\end{exercisesolution}

\begin{exercisesolution}[Solution to \cref{ex:box_normal_to_a_graph_in_r4}]
% Generator: Claude Opus 5 (High)
Write points of $\mathbb{R}^4$ as $(x,x_4)$ with $x \in \mathbb{R}^3$, and present the graph as a
level set: it is $G^{-1}(\{0\})$ for
\[G(x,x_4) := x_4 - \phi(x), \qquad
  \nabla G(x,x_4) = \begin{pmatrix} -\nabla\phi(x) \\ 1 \end{pmatrix} \in \mathbb{R}^4 .\]
This gradient never vanishes, its last entry being $1$, so
\cref{prop:gradient_orthogonality_level_sets} applies: the graph is a $3$-dimensional submanifold
of $\mathbb{R}^4$ and its normal space at $X$ is the line spanned by $\nabla G(X)$. Normalising,
\[N(X) = \frac{1}{\sqrt{1+\lvert\nabla\phi(x)\rvert^2}}
  \begin{pmatrix} -\partial_1\phi(x) \\ -\partial_2\phi(x) \\ -\partial_3\phi(x) \\ 1
  \end{pmatrix},\]
determined up to sign; this is the choice pointing in the direction of increasing $x_4$.

The formula is worth recognising rather than deriving. For $\phi : \mathbb{R}^2 \to \mathbb{R}$
it is the familiar unit normal $(-\partial_1\phi,-\partial_2\phi,1)/\sqrt{1+\lvert\nabla\phi
\rvert^2}$ of a surface in $\mathbb{R}^3$, and for $\phi : \mathbb{R} \to \mathbb{R}$ it is
$(-\phi',1)/\sqrt{1+\phi'^2}$, perpendicular to the tangent $(1,\phi')$ as it should be. Nothing
about the derivation used $n = 3$; presenting a graph as a level set of $x_{n+1} - \phi(x)$ works
in every dimension, which is the point of setting the question one dimension past the picture.

\begin{ainote}[A slip in the official key]
The key gives two displays for this answer. The second, $(-\nabla\phi(x),1)\transp$ divided by
$\sqrt{1+\lvert\nabla\phi(x)\rvert^2}$, is the formula above. The first is not: it lists a fourth
component $-\partial_4\phi(x)$ and normalises by
$\sqrt{1+\partial_1\phi(x)^2+\partial_2\phi(x)^2+\partial_3\phi(x)^2+\partial_4\phi(x)^2}$. Since
$\phi$ is a function of three variables there is no $\partial_4\phi$, and a vector with those
four components plus the entry $1$ would have five. Only the second display is correct.
\end{ainote}
\end{exercisesolution}


